\documentclass[12pt,a4paper]{article}

\usepackage[T2A]{fontenc}
\usepackage[utf8]{inputenc}
\usepackage[russian]{babel}
\usepackage[a4paper,margin=2.5cm]{geometry}
\usepackage{setspace}
\usepackage{parskip}

\setstretch{1.1}

\title{Текст доклада к презентации}
\author{}
\date{}

\begin{document}

\maketitle

\section*{Слайд 1. Титульный}

Здравствуйте. Тема моей выпускной квалификационной работы представлена на
слайде. Мой научный руководитель - Медведев Николай Викторович. Презентацию
также можно посмотреть по QR коду, даже на телефонах.

\section*{Слайд 2. Актуальность}

До первого клика, после заголовка:

Актуальность работы для меня во многом связана с практическим опытом.
Я работаю в VK Tech над СУБД Tarantool.

1-й клик — «Распределённые системы трудно проектировать и отлаживать.»

В рамках работ по синхронной репликации, которая реализована с
использованием алгоритма консенсуса Raft, был обнаружен критически
важный дефект. Сложность его анализа и исправления оказалась очень
высокой, потому что поведение такой системы зависит от редких сочетаний
событий, сетевых условий и состояний узлов.

2-й клик — «Ошибки в таких системах критичны для реальных сервисов.»

Для реального сервиса такие ошибки особенно опасны, поскольку они могут
приводить к нарушению согласованности данных, нестабильной работе
репликации и сложным сбоям, которые трудно воспроизвести обычным
тестированием.

3-й клик — «Формальная модель не гарантирует корректность кода.»

Для анализа проблемы было решено описать отдельные части Tarantool в
виде TLA+ спецификации, воспроизвести проблему на уровне формальной
модели, найти корректное решение сначала в спецификации, а затем
перенести его в исходный код. Но на практике возникла важная
трудность: даже если решение найдено в TLA+, это ещё не означает, что
реализация в коде действительно ему соответствует.

4-й клик — «Нужна методика сопоставления TLA+ спецификации и реализации.»

Именно это и привело к необходимости разработки методики, которая
позволяла бы систематически сопоставлять поведение программы с её
формальной моделью и проверять, что решение, полученное в спецификации,
корректно перенесено в реализацию.

\section*{Слайд 3. Цель и задачи}

Целью работы является разработка методики верификации соответствия
TLA+ спецификации и исходного кода, а также её проверка на
распределённой вычислительной системе.

Для достижения этой цели были поставлены следующие задачи.

Во-первых, необходимо было спроектировать и реализовать
распределённую систему, на которой можно практически проверить
предлагаемый подход.

Во-вторых, для этой системы требовалось разработать формальную
TLA+ спецификацию, которая выступает эталонной моделью поведения.

В-третьих, нужно было разработать и интегрировать систему
верификации, включающую инструментирование кода, построение трасс
исполнения и их последующую проверку.

И, наконец, в работе были рассмотрены экономические и правовые
аспекты разработки.

\section*{Слайд 4. Методика верификации}

На данном слайде представлена общая схема предлагаемой методики.
На входе имеются два основных артефакта.

Клик.

С одной стороны, это исходный код программной реализации.

Клик.

С другой стороны, это формальная спецификация на языке TLA+.

Клик.

Реализация инструментируется, и в процессе её выполнения формируются
трассы, содержащие логические шаги программы.

Клик.

Далее на основе исходной TLA+ спецификации строится спецификация
трассы, то есть ограниченная версия модели, которая разрешает только
те переходы, которые согласуются с наблюдаемой трассой.

Клик.

После этого средство TLC проверяет, существует ли допустимое поведение
модели, совместимое с данной трассой.

Если такое поведение существует, наблюдаемое исполнение считается
допустимым относительно спецификации. Если нет, значит между
реализацией и моделью обнаружено расхождение.

\section*{Слайд 5.1. Система распределённых вычислений}

Возвращаясь к общей схеме, сначала рассмотрим часть, связанную с
реализацией.

Здесь нас интересует объект, в который встраивается система
верификации, а именно разработанная распределённая система
вычислений.

\section*{Слайд 5.2. Система распределённых вычислений}

«На данном слайде представлена общая архитектура разработанной
распределённой системы вычислений.

Система состоит из трёх основных частей:
пользовательского клиента, кластера управляющих узлов Raft и
вычислительных узлов.

Центральным элементом является кластер Raft. Именно он поддерживает
согласованное состояние системы, принимает пользовательские запросы,
фиксирует изменения в реплицируемом журнале и координирует выполнение
задач.

Внутри кластера несколько Raft-узлов. Один из них в текущий момент
является лидером, а остальные поддерживают его состояние за счёт
репликации журнала.

Справа показаны вычислительные узлы. Они не участвуют в консенсусе и не
изменяют глобальное состояние напрямую. Их задача — выполнять полезную
нагрузку, которую им назначает кластер.

Слева расположен пользовательский клиент. Он взаимодействует не с
отдельными вычислительными узлами, а только с кластером Raft. Это
позволяет все изменения глобального состояния проводить через единый
согласованный механизм.

Таким образом, архитектура разделяет контур управления и контур
вычислений: кластер Raft отвечает за согласованность и координацию, а
вычислительные узлы — за непосредственное выполнение задач».

\section*{Слайд 5.3. Система распределённых вычислений}

Было проведено ручное и нагрузочное тестирование системы, с результатами
которого вы можете ознакомиться в отчете ВКР.

На слайде например приведен фрагмент узла лидера при нагрузочном тестировании
на большое количество клиентов. Мы видим монотонное увеличение индекса коммента
каждой новой записи лога.

Нас в данной системе интересует именно узел Raft, его корректность мы будем
сравнивать со спецификацей.

\section*{Слайд 6.1. Спецификация --- возврат к схеме}

Следующий компонент схемы --- это формальная спецификация.

Она играет роль эталона, с которым затем сопоставляется наблюдаемое
поведение программы.

\section*{Слайд 6.2. TLA+ спецификация алгоритма Raft}

Для управляющего кластера была разработана TLA+ спецификация алгоритма
консенсуса Raft. Эта спецификация описывает поведение кластера как систему
состояний и переходов.

Клик.

В ней формализованы основные переменные состояния каждого
узла: текущий терм, роль узла — последователь, кандидат или лидер,
журнал записей, индекс коммита.

Клик.

Кроме локального состояния узлов в спецификации моделируется и среда
взаимодействия — то есть сообщения, которые циркулируют между узлами.
За счёт этого можно описывать не только внутреннее состояние каждого
сервера, но и сам процесс распределённого обмена сообщениями.


В спецификации заданы основные действия алгоритма Raft.

Клик

Это, прежде всего, запуск выборов лидера,

Клик

отправка и обработка запросов голосов,

Клик

отправка и обработка данных и пингов

Клик

Именно совокупность этих действий задаёт допустимые переходы модели из
одного состояния в другое.

Отдельно в спецификации сформулированы инварианты, то есть свойства,
которые должны сохраняться при любых допустимых переходах.

Клик

Первый важный инвариант связан с согласованностью журналов. Его смысл
состоит в том, что если у двух серверов некоторый индекс уже входит в
зафиксированную часть журнала, то записи по всем индексам этого общего
префикса должны совпадать. Иначе говоря, после фиксации команда не
может иметь разные значения на разных узлах кластера.

Второй важный инвариант связан с сохранностью уже подтверждённых
значений. Если некоторое значение уже считается подтверждённым, то
любой актуальный лидер обязан содержать это значение в своём журнале.
Это исключает ситуацию, при которой после смены лидера система теряет
данные, о фиксации которых уже было сообщено клиенту.

Таким образом, спецификация описывает не только структуру состояния и
механизм работы Raft, но и формально задаёт те свойства безопасности,
которые обязаны выполняться при любом сценарии исполнения.

\section*{Слайд 6.1. TLA+ спецификация алгоритма Raft}

Эти инварианты были проверены с помощью TCL, было исследовано более 36 тысяч
различных состояний, нарушений не обнаружено, а потому можно считать
специфкацию рабочей.

Именно эта спецификация далее используется как эталонная модель:
с одной стороны, она может быть проверена средствами TLC, а с другой —
она служит формальной основой для последующего сопоставления с
наблюдаемыми трассами исполнения реальной реализации

\section*{Слайд 7.1. Методика верификации --- возврат к схеме}

Следующий важный элемент методики --- это трассы исполнения.

Чтобы проверить реализацию, нужно сначала получить из неё
последовательность логических шагов в специальном формате.

\section*{Слайд 7.2. Библиотека инструментирования и формат трассы}

Для этого была разработана библиотека инструментирования.

Её работа основана на двух основных операциях.

Клик.

Первая операция --- фиксация изменений логических переменных модели.
В коде это делается через вызов \texttt{notify\_change}.

Клик.

Вторая операция --- вызов \texttt{log} в точке линеаризации, то есть в
момент, когда завершён один логически атомарный шаг программы.
В результате каждая запись трассы содержит временную метку, имя
события, аргументы события и набор обновлений переменных.

Клик.

Для согласования порядка событий в библиотеке используется механизм
часов. В реализации предусмотрены два варианта: часы в памяти,
синхронизируемые через мьютексы, и файловые часы, которые могут
использоваться несколькими процессами на одном компьютере. Это
позволяет назначать записям монотонно возрастающие временные метки и
сохранять корректный порядок логических шагов как внутри одного
процесса, так и между несколькими процессами

Клик.

Таким образом, одна запись трассы соответствует одному логическому
наблюдению, которое затем сопоставляется с переходом TLA+ модели.
Если трассы формируются несколькими процессами, они объединяются в
одну глобальную трассу. Пример записи трассы приведен на слайде.

\section*{Слайд 8.1. Методика верификации --- возврат к схеме}

Теперь перейдём к следующему ключевому элементу методики ---
спецификации трассы. Именно здесь происходит проверка соответствия
спецификации и реализации, а значит настало время сформулировать
математическую постановку задачи.

\section*{Слайд 8.1. Математическая постановка задачи}

На этом слайде приведена математическая постановка задачи.

Здесь M — это формальная модель программы в виде системы
переходов, то есть множество её состояний, начальных состояний и
допустимых переходов между ними.

tau — это трасса исполнения программы, то есть последовательность
наблюдаемых шагов, полученных из реального выполнения реализации.

C — это предикат согласованности шага. Он показывает, можно ли
считать очередное наблюдение трассы согласованным с некоторым переходом
формальной модели.

Далее вводится множество K(M,tau), то есть множество
всех допустимых длин префиксов трассы, которые ещё можно согласовать с
моделью.

Тогда задача состоит в том, чтобы найти такое значение k, при
котором показатель качества согласованности максимален.

Говоря проще, мы ищем максимальную длину префикса трассы, которую ещё
можно объяснить корректным поведением модели.

Если таким префиксом является вся трасса, то реализация полностью
согласуется со спецификацией. Если нет, значит начиная с некоторого
шага возникает расхождение между кодом и моделью.

\section*{Слайд 8.2. Решение поставленной задачи}

Решение построено через последовательность множеств $D_i$.

Множество $D_0$ --- это множество начальных состояний модели.

Далее для каждого следующего шага трассы строится множество $D_i$,
состоящее из тех состояний, в которые можно перейти из предыдущего
множества $D_{i-1}$ допустимым переходом модели и которые
согласуются с текущим наблюдением трассы.

Если на некотором шаге оказывается, что $D_i$ пусто, это означает, что
дальнейшее согласование невозможно, и алгоритм останавливается.

Максимальное значение $k^*$, до которого множества остаются
непустыми, и есть длина максимального согласованного префикса.

Таким образом, предложенный метод позволяет не только ответить,
совместима трасса целиком или нет, но и определить первый шаг, на
котором возникает расхождение между реализацией и спецификацией.

\section*{Слайд 8.3. Спецификация трассы}

Давайте рассмотрим, каким образом реализуется решение математической задачи в
коде TLA+.

Модель читает переданную ей трассу программы и сохраняет ее в переменную Trace.

Клик

Сначала TLC запускает не базовую спецификацию Raft, а ограничение над
спецификацей TraceSpec.

Клик

Далее в TraceSpec строится начальное состояние TraceInit.
Оно состоит из двух частей:
во-первых, индекс текущей записи трассы l устанавливается в 1,
а во-вторых, выполняется BaseInit, то есть обычное начальное
состояние спецификации Raft.

Иными словами, в начале проверки модель находится в своём
корректном начальном состоянии и готова обрабатывать первую
запись трассы.

Клик

Здесь основной интерес представляет TraceNext.
В конфигурации он переопределён как RaftTraceNext.
RaftTraceNext — это дизъюнкция всех допустимых действий Raft,
которые могут соответствовать текущей записи трассы.

Клик

Каждый такой предикат работает одинаково.
Сначала он через IsEvent(...) проверяет, что имя события в
текущей строке трассы совпадает с ожидаемым.
Если в трассе заданы аргументы события, например идентификатор
узла, то выбирается конкретный экземпляр действия модели.
Если аргументы не заданы, допускается любой подходящий экземпляр.

Клик

Внутри IsEvent(...) происходит ещё две важные вещи.
Во-первых, индекс трассы увеличивается: l' = l + 1,
то есть модель переходит к следующей записи трассы.
Во-вторых, вызывается UpdateVariables(logline), а в нашем случае
это RaftUpdateVariables.

Клик

RaftUpdateVariables берёт текущую запись трассы и, если в ней
зафиксированы изменения переменных, накладывает на них точные
ограничения

Клик

Если базовое действие Raft пытается обновить такую переменную иначе, чем это
зафиксировано в трассе, то данный переход становится недопустимым. Для
переменных, которые в трассе не указаны, значение выбирается не произвольно, а
только из тех значений, которые допускает выбранное действие модели.

Клик

Таким образом, один шаг проверки устроен так:
TLC берёт текущую запись трассы, ищет такое действие Raft,
которое совпадает по имени события, подходит по аргументам и даёт
нужные обновления переменных.
Если такой шаг существует, модель переходит дальше.
Если нет, дальнейшее согласование становится невозможным, выводится номер
записи трассы и сама запись трассы.

Именно так реализуется решение математической постановки:
на каждом шаге остаются только те состояния, которые достижимы
из предыдущих и согласуются с очередным наблюдением трассы.
То есть TLC неявно строит те самые множества согласованных
состояний по префиксам трассы $D_i$.

\section*{Слайд 9. Тестирование системы верификации}

В ходе тестирования системы верификации удалось выявить не просто
формальные расхождения, а реальные ошибки реализации.

На демонстрационном примере Two-Phase Commit была обнаружена критичная
ошибка: из-за повторной отправки сообщения Prepared координатор мог
принять решение Commit раньше, чем получит подтверждение от всех
участников. Это критично, потому что нарушается атомарность транзакции:
часть узлов может считать транзакцию завершённой, а часть — нет.

На реализации Raft также были найдены два несоответствия спецификации:
неверное начальное значение currentTerm и отправка RequestVote
самому себе. Первая ошибка нарушает согласованность логики термов уже с
первого шага исполнения. Вторая опасна тем, что узел может фактически
учесть собственный голос дважды и преждевременно получить лидерство.

Клик.

Важно, что все эти расхождения были локализованы по конкретным шагам
трассы. То есть методика позволяет не только сказать, что реализация не
соответствует спецификации, но и показать, где именно возникает
проблема». Пример несоответсвия запроса голосов в Raft приведен на слайде.

\section*{Слайд 10. Достоинства и недостатки}

На этом слайде приведены основные достоинства и ограничения
предложенной методики.

К её достоинствам можно отнести, во-первых, то, что она позволяет
сопоставлять не абстрактную модель саму по себе, а именно реальное
исполнение программы с TLA+ спецификацией. Во-вторых, методика
позволяет локализовать проблему по конкретному шагу трассы, то есть
не просто зафиксировать наличие расхождения, а указать, на каком
наблюдении оно возникает. В-третьих, для её применения не требуется
обязательно логировать все переменные модели. Можно фиксировать только
существенную часть состояния, а недостающие компоненты будут
восстанавливаться в рамках тех переходов, которые допускает
спецификация.

Однако у подхода есть и ограничения. Важное ограничение связано с полнотой
трассы. Если в записи трассы содержится слишком мало информации, например
зафиксировано только имя события, но почти не зафиксированы изменения
состояния, то множество допустимых интерпретаций такого шага становится слишком
широким. В этом случае TLC действительно может найти некоторое согласованное
поведение модели, но фактическое поведение реализации при этом будет
восстановлено неоднозначно.

Поэтому практическая точность методики зависит не только от
корректности самой спецификации, но и от качества инструментирования
кода: от того, насколько удачно выбраны точки линеаризации, какие
переменные логируются и насколько подробно описываются наблюдаемые
переходы. Иными словами, методика работает тем точнее, чем более
информативную трассу удаётся получить из реализации

\section*{Слайд 11. Организационно-экономическая часть}

На данном слайде представлены итоговые технико-экономические
показатели проекта. Общая трудоёмкость разработки составила
436 человеко-часов, что соответствует 54,5 человеко-дням, а
календарная продолжительность проекта — 43 рабочих дня. Наиболее
трудоёмким этапом стала разработка библиотеки трассировки.

Суммарные затраты на реализацию проекта оцениваются примерно в
3,98 миллиона рублей. При принятой цене одной копии продукта и
планируемом объёме продаж расчёты показывают положительные
экономические результаты: проект достигает точки безубыточности
примерно за 13–15 месяцев, а срок окупаемости инвестиций составляет
около 8 месяцев.

Таким образом, выполненная оценка показывает, что разработка является
не только технически значимой, но и экономически целесообразной

\section*{Слайд 12. Организационно-правовое обеспечение ИБ}

«В правовой части были рассмотрены основные нормативно-правовые акты,
которые относятся к теме работы.

Прежде всего это Конституция Российской Федерации, которая закрепляет
базовые права в информационной сфере, в том числе право на защиту
частной жизни и конфиденциальность информации.

Далее был рассмотрен Федеральный закон №149-ФЗ, который регулирует
отношения в области информации, информационных технологий и её защиты.
Для данной работы он важен, поскольку методика связана со сбором,
обработкой и анализом трасс исполнения.

Также был учтён Указ Президента №646, утвердивший Доктрину
информационной безопасности Российской Федерации. Он задаёт общий
стратегический контекст, в котором повышение надёжности и корректности
программных систем рассматривается как значимая задача.

Кроме того, были проанализированы ГОСТ Р 56939–2024 и приказ ФСТЭК
России №240, которые относятся к безопасной разработке программного
обеспечения и процессам её подтверждения.

В целом проведённый анализ показал, что разработанная методика
соответствует действующим нормативно-правовым требованиям Российской
Федерации и может рассматриваться как допустимый инструмент безопасной
разработки программного обеспечения».

\section*{Слайд 13. Заключение и перспективы}

Подводя итог, можно выделить следующие результаты работы.

Была разработана методика верификации соответствия TLA+ спецификации и
исходного кода.

Была сформулирована математическая постановка задачи и предложен
алгоритм её решения через поиск максимального согласованного
префикса трассы.

Были реализованы распределённая система вычислений на основе Raft,
формальная TLA+ спецификация и система верификации, включающая
библиотеку инструментирования и механизм проверки трасс.

Проведённая апробация показала, что предложенный подход позволяет
выявлять реальные логические ошибки реализации и применять формальные
методы не только на уровне спецификации, но и на уровне кода.

В качестве дальнейшего развития работы планируется расширение
библиотеки инструментирования на новые языки программирования,
прежде всего на Lua, а также интеграция предложенной методики в
СУБД Tarantool, где Lua широко используется для прикладной логики.

Спасибо за внимание.

\end{document}
```
